GBFPlus supports authorized tracing of cross-domain data transfers through
public ledger attestations. Its published description treats the adjacent-domain
identifier as protected metadata. We show that an allowed one-insertion
attestation reveals that identifier exactly. Public occupancy identifies all
shares of the inserted word; their exclusive-or (XOR) reconstructs the word,
including its unmasked identifier suffix. This holds for every construction
realization and requires neither a query token nor a candidate data identifier.
A public occupancy test also identifies some larger batches in which the same
operation reveals an aggregate XOR, rather than individual identifiers.
We prove these statements and reproduce them with finite exhaustive checks and
stored share arrays. The disclosure concerns adjacent-domain metadata, distinct
from the source paper's data-identifier privacy game and signer anonymity.
The stated metadata-confidentiality claim therefore needs to account for allowed
partial batches and the complete public payload view.
